% GENERATED FILE. Edit canonical lessons, then run npm run generate:books.
% canonical-source-hash: fnv1a64:e712cc17887fd6d0
% canonical-lessons: SA-C247-banah, SA-C247-khadgah, SA-C247-kavacam, SA-C247-dhvajah, SA-C247-patahah

\chapter{Arrow, Sword, Armour, Flag, Drum}
\label{ch:sa-arrow-sword-armour-flag-drum}
\hlchaptermodality{247}

\begin{chapteropening}
\textbf{By the end of this chapter:} \emph{I can say an arrow, a sword, armour, a flag and a drum.}

Complete the last lesson of chapter 247: I can say an arrow, a sword, armour, a flag and a drum.
\end{chapteropening}

\section[\texorpdfstring{\sk{बाणः} (bāṇaḥ)}{बाणः (bāṇaḥ)}]{\sk{बाणः} (bāṇaḥ) — an arrow}

\label{lesson:SA-C247-banah}



\begin{quote}
Before the new one: say the Sanskrit for wild honey, then the Sanskrit for a sweet dish.
\end{quote}

\subsection*{You'll want to know: \sk{बाणः}}
\textbf{\sk{बाणः}} — \emph{bāṇaḥ} — \textquotedblleft{}an arrow\textquotedblright{}.

\begin{grammarlens}[title={What you've built}]
One more word for this chapter.
\end{grammarlens}

\subsection*{Your turn}
\begin{itemize}
\raggedright
  \item \emph{Say it:} \emph{bāṇaḥ}
  \item \emph{Say it:} \emph{bāṇaḥ}, once more
  \item \emph{Say it:} say \emph{bāṇaḥ}
  \item \emph{Recall:} say the Sanskrit for wild honey, then the Sanskrit for a sweet dish, then say \emph{bāṇaḥ} again
\end{itemize}

\begin{tcolorbox}[breakable,colback=teal!4,colframe=teal!35!black,title={Before you move on}]
What does \sk{बाणः} mean? (\textquotedblleft{}An arrow\textquotedblright{}.) Say it once more.
\end{tcolorbox}
\section[\texorpdfstring{\sk{खड्गः} (khaḍgaḥ)}{खड्गः (khaḍgaḥ)}]{\sk{खड्गः} (khaḍgaḥ) — a sword}

\label{lesson:SA-C247-khadgah}



\begin{quote}
Before the new one: say the Sanskrit for a sweet dish, then the Sanskrit for an arrow.
\end{quote}

\subsection*{You'll want to know: \sk{खड्गः}}
\textbf{\sk{खड्गः}} — \emph{khaḍgaḥ} — \textquotedblleft{}a sword\textquotedblright{}.

\begin{grammarlens}[title={What you've built}]
One more word for this chapter.
\end{grammarlens}

\subsection*{Your turn}
\begin{itemize}
\raggedright
  \item \emph{Say it:} \emph{khaḍgaḥ}
  \item \emph{Say it:} \emph{khaḍgaḥ}, once more
  \item \emph{Say it:} say \emph{khaḍgaḥ}
  \item \emph{Recall:} say the Sanskrit for a sweet dish, then the Sanskrit for an arrow, then say \emph{khaḍgaḥ} again
\end{itemize}

\begin{tcolorbox}[breakable,colback=teal!4,colframe=teal!35!black,title={Before you move on}]
What does \sk{खड्गः} mean? (\textquotedblleft{}A sword\textquotedblright{}.) Say it once more.
\end{tcolorbox}
\section[\texorpdfstring{\sk{कवचम्} (kavacam)}{कवचम् (kavacam)}]{\sk{कवचम्} (kavacam) — armour}

\label{lesson:SA-C247-kavacam}



\begin{quote}
Before the new one: say the Sanskrit for an arrow, then the Sanskrit for a sword.
\end{quote}

\subsection*{You'll want to know: \sk{कवचम्}}
\textbf{\sk{कवचम्}} — \emph{kavacam} — \textquotedblleft{}armour\textquotedblright{}.

\begin{grammarlens}[title={What you've built}]
One more word for this chapter.
\end{grammarlens}

\subsection*{Your turn}
\begin{itemize}
\raggedright
  \item \emph{Say it:} \emph{kavacam}
  \item \emph{Say it:} \emph{kavacam}, once more
  \item \emph{Say it:} say \emph{kavacam}
  \item \emph{Recall:} say the Sanskrit for an arrow, then the Sanskrit for a sword, then say \emph{kavacam} again
\end{itemize}

\begin{tcolorbox}[breakable,colback=teal!4,colframe=teal!35!black,title={Before you move on}]
What does \sk{कवचम्} mean? (\textquotedblleft{}Armour\textquotedblright{}.) Say it once more.
\end{tcolorbox}
\section[\texorpdfstring{\sk{ध्वजः} (dhvajaḥ)}{ध्वजः (dhvajaḥ)}]{\sk{ध्वजः} (dhvajaḥ) — a flag}

\label{lesson:SA-C247-dhvajah}



\begin{quote}
Before the new one: say the Sanskrit for a sword, then the Sanskrit for armour.
\end{quote}

\subsection*{You'll want to know: \sk{ध्वजः}}
\textbf{\sk{ध्वजः}} — \emph{dhvajaḥ} — \textquotedblleft{}a flag\textquotedblright{}.

\begin{grammarlens}[title={What you've built}]
One more word for this chapter.
\end{grammarlens}

\subsection*{Your turn}
\begin{itemize}
\raggedright
  \item \emph{Say it:} \emph{dhvajaḥ}
  \item \emph{Say it:} \emph{dhvajaḥ}, once more
  \item \emph{Say it:} say \emph{dhvajaḥ}
  \item \emph{Recall:} say the Sanskrit for a sword, then the Sanskrit for armour, then say \emph{dhvajaḥ} again
\end{itemize}

\begin{tcolorbox}[breakable,colback=teal!4,colframe=teal!35!black,title={Before you move on}]
What does \sk{ध्वजः} mean? (\textquotedblleft{}A flag\textquotedblright{}.) Say it once more.
\end{tcolorbox}
\section[\texorpdfstring{\sk{पटहः} (paṭahaḥ)}{पटहः (paṭahaḥ)}]{\sk{पटहः} (paṭahaḥ) — a drum}

\label{lesson:SA-C247-patahah}



\begin{quote}
Before the new one: say the Sanskrit for armour, then the Sanskrit for a flag.
\end{quote}

\subsection*{You'll want to know: \sk{पटहः}}
\textbf{\sk{पटहः}} — \emph{paṭahaḥ} — \textquotedblleft{}a drum\textquotedblright{}.

\begin{grammarlens}[title={What you've built}]
One more word for this chapter.
\end{grammarlens}

\subsection*{Your turn}
\begin{itemize}
\raggedright
  \item \emph{Say it:} \emph{paṭahaḥ}
  \item \emph{Say it:} \emph{paṭahaḥ}, once more
  \item \emph{Say it:} say \emph{paṭahaḥ}
  \item \emph{Recall:} say the Sanskrit for armour, then the Sanskrit for a flag, then say \emph{paṭahaḥ} again
\end{itemize}

\begin{tcolorbox}[breakable,colback=teal!4,colframe=teal!35!black,title={Before you move on}]
What does \sk{पटहः} mean? (\textquotedblleft{}A drum\textquotedblright{}.) Say it once more.
\end{tcolorbox}
